\section{Adaptive Evidence Execution}
\label{sec:execution}

% AE-FIG — System architecture (static schematic; same tabular style as Fig3)
\begin{figure}[t]
\centering
\small
\begin{tabular}{@{}c@{}}
\framebox{\parbox{0.92\columnwidth}{\centering
\textbf{Question} $\to$ \textbf{Slot Compiler} (LLM) \\
$\downarrow$ typed evidence requirements + dependency edges}} \\
$\downarrow$ \\
\framebox{\parbox{0.92\columnwidth}{\centering
\textbf{Declarative Evidence Algebra} (Section~\ref{sec:algebra}) \\
states, types, operators, logical/physical plan split}} \\
$\downarrow$ \\
\framebox{\parbox{0.92\columnwidth}{\centering
\textbf{Requirement-Aware Allocator} (Section~\ref{sec:optimizer}) \\
importance $\tau{=}2d{-}1$, budgeted allocation, Pareto search}} \\
$\downarrow$ \\
\framebox{\parbox{0.92\columnwidth}{\centering
\textbf{Adaptive Executor} (Section~\ref{sec:execution}) \\
retrieval $\to$ join $\to$ materialize $\to$ generate, under hard budget}} \\
$\downarrow$ \\
\framebox{\parbox{0.92\columnwidth}{\centering
\textbf{Structured Answer Contract} $\to$ evidence-typed answer + provenance}} \\
\end{tabular}
\caption{End-to-end SlotRAG pipeline. The declarative evidence program is the
interface between the compiler and the requirement-aware allocator; the
adaptive executor realizes the chosen physical plan under an explicit
retrieval budget. This is a static schematic of the system, not a result
plot.}
\label{fig:architecture}
\end{figure}

Figure~\ref{fig:architecture} shows the end-to-end pipeline. A frozen logical
plan and a chosen physical plan are executed by the adaptive executor.

% AE-P1 — Deterministic adaptive execution (honest rewrite of G4)
Under our design, execution is \emph{deterministic}: given the same
logical plan, the same budget, and the same corpus, the physical-plan selection
is a pure function of requirement importance and remaining budget. Runtime
observations---a binding that is found or missing, a candidate yield that
differs from the estimate---are recorded as evidence state and telemetry, but
under the current chain-only execution topology they do not re-allocate the
physical plan. We state this boundary explicitly. A strict chain admits exactly
one legal execution order, so runtime re-optimization over alternative orders
has nothing to decide; implementing it on this topology would be dead code. We
therefore treat runtime re-optimization as future work under a branching
execution model, and the contribution of this section is the deterministic
execution semantics, budget enforcement, and provenance that make execution
auditable rather than the re-optimization machinery.

% AE-P2 — Executor state transition
Each physical operator transforms the evidence state $S_t$ into $S_{t+1}$ while
recording realized cost, candidate yield, verification outcome, bindings, and
provenance. The state transition is governed by the algebra of
Section~\ref{sec:algebra}: \code{apply\_observation} advances each requirement
along the monotone lattice \code{unresolved} $\rightarrow$ \code{partial}
$\rightarrow$ \code{satisfied}, and never regresses. The evidence state is a
first-class runtime snapshot (\code{EvidenceState}) containing the requirement
statuses, the bound evidence, the current bindings, and the realized budget
usage. It is derived from the immutable per-slot execution traces
(\code{derive\_evidence\_state}), so every status transition is auditable
against the trace that produced it. This is the mechanism-level diagnostic that
distinguishes \emph{acquisition} failure (nothing bound) from \emph{selection}
failure (evidence retrieved but not sufficient) from \emph{integration} failure
(evidence sufficient but the answer generator failed to use it).

% AE-P3 — Re-optimization triggers → future work
Event-driven re-optimization over alternative physical plans is not part of the
current system. We identify the trigger design---observed selectivity or cost
deviation beyond a threshold, a requirement changing status, a new binding
enabling a previously infeasible operator, or a budget shock invalidating the
current plan---as the specification for future work under a branching execution
model, where multiple legal orders make re-optimization non-vacuous.

% AE-P4 — Partial-plan preservation → future work
In a branching execution model, re-optimization should preserve valid
materialized evidence and completed bindings, reconsidering only the unexecuted
suffix of the physical plan. We leave this to future work as well: on the
current single-order topology the suffix is trivial, and we prefer to report
the deterministic behavior we actually implemented over machinery that would
not change any executed plan.

% AE-P5 — Provenance and diagnostics (contribution)
Provenance is a system property, not a debugging aid. Every materialized
evidence record carries its source id; every slot execution trace records its
retrieval candidates, binding contexts, sufficiency verdict, and materialization
outcome; and the derived evidence state records, for each non-satisfied
requirement, the reason it is unsatisfied (slot never materialized, no binding
context attempted, calibrator verdict of insufficient evidence, or partial
coverage). This lets a query be traced end-to-end from answer to the physical
actions and sources that produced it, and lets failure attribution be
operator-level rather than a single opaque ``wrong answer'' label. The
distinction between retrieved, materialized, and packed evidence
(Section~\ref{sec:algebra}) is what makes such attribution sound: a stage is
blamed only when the evidence it was supposed to produce is genuinely absent
from the downstream state.

% AE-P6 — Budget enforcement and stopping (contribution)
Budget enforcement is deterministic. The executor threads the global
retrieval-call budget into the optimizer's search (\code{retrieval\_budget}),
so the physical plan is selected under the same budget that bounds execution.
It reserves calls for unresolved high-priority requirements, rejects plans that
would violate hard constraints, and terminates when requirements are
sufficiently satisfied or no admissible plan can improve utility within the
remaining budget. This is not a learned stopping controller; it is a
deterministic, auditable rule. Because the whole path---plan selection, budget
threading, execution, state derivation, and stopping---is deterministic given a
frozen logical plan, the system admits exact reproduction and per-question
trace audits, which we exploit in the evaluation (Section~\ref{sec:experiments}).
